\chapter{Introduction to AULinux}
\label{chap:intro}

\noindent
\emph{AULinux}---``Absolutely Unique Linux''---is a custom Linux distribution
built from the ground up around a single, deliberately unfashionable idea: that
no one programming language is the right tool for every layer of an operating
system. Where most distributions inherit a userland written almost entirely in
C, AULinux assembles its userspace from three languages, each chosen for the
layer it serves best. C drives the metal-adjacent code---the init system and the
loadable kernel modules. Go provides a statically linked, memory-safe core
userland: the coreutils replacements and the interactive shell. Rust implements
the package manager, where the borrow checker turns filesystem and state
manipulation from a minefield into a tractable engineering problem.

This book is a technical deep-dive into how those pieces actually fit together.
Every chapter is grounded in the real source tree---the C of \file{init/src/init.c},
the kernel modules under \file{kernel/src/}, the Go in \file{utils/} and
\file{shell/}, and the Rust in \file{pkg-manager/src/}. Where a feature is
fully implemented, we walk its control flow and read its code. Where a feature
is aspirational or stubbed, we say so plainly. The goal is not a marketing tour
but an honest architectural map a serious engineer can navigate and extend.

\begin{infobox}[The thesis in one sentence]
A modern OS spans radically different abstraction layers; AULinux matches each
layer to a language whose guarantees fit that layer's failure modes---C for raw
control, Go for safe and self-contained userland binaries, Rust for stateful
system management that must never corrupt the machine.
\end{infobox}

\section{The Polyglot Triad}

The three languages are not used interchangeably. Each owns a tier of the system
and is held responsible for a distinct class of correctness.

\begin{description}[leftmargin=1.4em, style=nextline]
  \item[\textcolor{aukernel}{C --- the privileged tier}]
    The first userspace process (\code{au-init}, PID~1) and the in-kernel
    modules (\code{aulinux\_core}, \code{aulinux\_procmon}, \code{aulinux\_sysctl})
    are written in C. At this tier there is no runtime to lean on: PID~1 runs
    before any dynamic linker is available, and kernel modules execute in
    ring~0 where a stray pointer panics the whole machine. C's transparency---its
    near one-to-one mapping to syscalls and memory---is exactly what this tier
    needs.
  \item[\textcolor{auteal}{Go --- the userland tier}]
    The coreutils replacements (\code{ls}, \code{cat}, \code{cp}, \code{ps},
    \dots) and the shell (\code{aush}) are written in Go and compiled, BusyBox
    style, into one \emph{multicall} binary called \code{auutils}. Go's effortless
    static linking yields binaries with zero shared-library dependencies---no
    missing \file{libc.so.6}, no \code{ld-linux} at all---while its garbage
    collector and bounds checking eliminate the buffer overflows that have
    haunted C coreutils for decades.
  \item[\textcolor{aurust}{Rust --- the stateful-management tier}]
    The package manager \code{aupkg} is written in Rust. Installing and removing
    software means mutating the live filesystem and a persistent database; a
    crash mid-transaction can brick a system. Rust's ownership model and
    \code{Result}-based error handling make resource cleanup and error
    propagation structural rather than aspirational.
\end{description}

\begin{table}[h]
\centering
\renewcommand{\arraystretch}{1.3}
\begin{tabular}{>{\raggedright\arraybackslash}p{2.6cm} l l >{\raggedright\arraybackslash}p{4.2cm}}
\toprule
\textbf{Component} & \textbf{Language} & \textbf{Artifact} & \textbf{Primary responsibility} \\
\midrule
Init system    & C    & \code{au-init} (static)   & PID~1: mount, supervise, reap \\
Kernel modules & C    & \code{*.ko}               & Telemetry, sysctls, procmon \\
Core utilities & Go   & \code{auutils} (multicall) & Coreutils replacements \\
Shell          & Go   & \code{auutils} (\code{aush}) & Interactive REPL, job control \\
Package manager& Rust & \code{aupkg}              & Install / remove / query packages \\
Build pipeline & Bash & \code{scripts/*.sh}       & Compile, assemble rootfs, boot \\
\bottomrule
\end{tabular}
\caption{The AULinux component map and the language that owns each tier.}
\end{table}

\section{System Architecture at a Glance}

AULinux boots a stock, unmodified Linux kernel and layers its identity entirely
in the surrounding software. The runtime stack, from power-on upward, looks like
this:

\begin{enumerate}
  \item \textbf{GRUB} loads the Linux kernel and an \code{initramfs} containing
        the AULinux root filesystem.
  \item The \textbf{kernel} initialises hardware, then hands control to
        \code{/sbin/init}---the statically linked C program \code{au-init}.
  \item \textbf{\code{au-init} (PID~1)} mounts the pseudo-filesystems, brings up
        loopback networking, reads \file{/etc/aulinux/services.conf}, spawns the
        login shell and any declared services, and enters a supervision loop that
        reaps zombies and restarts crashed children.
  \item The \textbf{\code{aush}} shell---one personality of the \code{auutils}
        multicall binary---gives the user an interactive prompt, dispatching to
        the other coreutils personalities and to external programs.
  \item \textbf{\code{aupkg}}, the Rust package manager, manages software on top
        of all this, reading and writing its JSON databases under
        \file{/var/lib/aupkg}.
  \item Optional \textbf{kernel modules} (\code{aulinux\_core} and friends) can be
        \code{insmod}'d to expose AULinux-specific telemetry under \file{/proc}.
\end{enumerate}

\begin{notebox}[Standard kernel, custom everything else]
AULinux does not fork or patch the upstream Linux kernel. Its kernel
customisation lives entirely in out-of-tree loadable modules
(Chapter~\ref{chap:kernel} discusses why). This keeps AULinux compatible with
mainline kernel security updates while still allowing bespoke instrumentation.
\end{notebox}

\section{Design Philosophy}

Three principles recur throughout the system and are worth stating up front,
because they explain decisions you will meet again in every chapter.

\subsection{Static linking everywhere it counts}
\code{au-init} is linked with \code{-static} so PID~1 can run before any dynamic
linker exists; the Go binaries are built with \code{CGO\_ENABLED=0} and
\code{-extldflags '-static'} so a single ELF file works anywhere in the image;
\code{aupkg} is compiled with \code{crt-static}. The payoff is relocatability and
the elimination of an entire category of ``missing \file{.so}'' boot failures.
The cost---larger binaries---is clawed back by the multicall design and optional
UPX compression (Chapter~\ref{chap:build}).

\subsection{Right-sized footprint via the multicall pattern}
Rather than ship a dozen separate Go executables, AULinux compiles every utility
and the shell into one \code{auutils} binary and populates \file{/bin} with
symlinks pointing at it. The program inspects \code{argv[0]} to decide which
personality to become---the same trick BusyBox uses---so the on-disk cost of
fifteen tools is one binary plus a handful of symlinks.

\subsection{Honesty about maturity}
AULinux is a real, bootable system, but it is also a young one. Some subsystems
are complete (the core kernel telemetry module, the shell's pipeline engine, the
package database's atomic writes); others are partial (the process-monitor
module's fork probe, \code{aupkg}'s update/upgrade commands). This book flags the
difference at every turn using the callouts you have already seen, so that the
map matches the territory.

\begin{warnbox}[Read the noteboxes]
Throughout the book, blue \emph{Note} boxes mark caveats and clarifications, and
amber \emph{Caution} boxes mark genuine limitations, bugs, or stubbed features
found in the present source tree. They are not decoration---they are where the
gap between intent and implementation is recorded.
\end{warnbox}

\section{How This Book Is Organised}

The remaining chapters follow the system from the bottom up:

\begin{itemize}
  \item \textbf{Chapter~\ref{chap:init} --- Bootstrapping and Initialization:}
        the boot chain and a line-by-line tour of \code{au-init}, PID~1.
  \item \textbf{Chapter~\ref{chap:kernel} --- Dynamic Kernel Extensions:}
        the three loadable kernel modules and the \file{/proc} surfaces they
        expose.
  \item \textbf{Chapter~\ref{chap:userland} --- The Go-Powered Userland:}
        the \code{auutils} multicall binary and the \code{aush} shell.
  \item \textbf{Chapter~\ref{chap:pkgmgr} --- The Rust Package Manager:}
        \code{aupkg}'s CLI, data model, and on-disk database.
  \item \textbf{Chapter~\ref{chap:build} --- Assembly and the Build Pipeline:}
        how the polyglot pieces are compiled, packed into a rootfs, and booted
        in QEMU.
\end{itemize}

\noindent
With the map drawn, we begin where the machine itself begins: at power-on, with
the handoff from firmware to kernel to the first process of userspace.
